% =====================================================================
% 06 — Training methods for EBMs.
% =====================================================================
% Planning (% comments only):
% For each: equation, what it does, where the partition function moves.
% Cover: MLE with MCMC, contrastive divergence, score matching,
% noise-contrastive estimation, denoising score matching,
% energy-based score matching, ratio matching, persistent CD, JEM.
% =====================================================================

\section{Training Energy-Based Models}
\label{sec:ebm_training}

The training methods for energy-based models can be organized by what they do with $\Ztheta$.
\Cref{tab:ebm_training} previews the taxonomy.
We then take each method in turn.

\begin{table}[ht]
  \centering
  \small
  \caption{Training methods for energy-based models, organized by their treatment of the partition function.
  ``Cost moved to'' indicates where the difficulty appears in practice.}
  \label{tab:ebm_training}
  \begin{tabular}{@{}lll@{}}
    \toprule
    Method & Treatment of $\Ztheta$ & Cost moved to \\
    \midrule
    MLE with sampling                   & estimated via MCMC samples           & MCMC mixing per gradient step \\
    Contrastive divergence              & truncated MCMC chains                & bias from short chains \\
    Persistent CD                       & long-lived chains across updates     & chain decorrelation \\
    Noise-contrastive estimation        & cast as binary classification        & quality of the noise distribution \\
    Score matching                      & avoids $\Ztheta$ entirely            & second-order derivatives of the energy \\
    Denoising score matching            & avoids $\Ztheta$ entirely            & noise schedule design \\
    Sliced score matching               & avoids $\Ztheta$ entirely            & random-projection variance \\
    Ratio matching                      & avoids $\Ztheta$ entirely (discrete) & per-coordinate ratio comparisons \\
    Stein discrepancy                   & avoids $\Ztheta$ entirely            & kernel choice \\
    Implicit differentiation / JEM      & shares $\Ztheta$ with a classifier   & inference for the gradient \\
    \bottomrule
  \end{tabular}
\end{table}

\subsection{Maximum likelihood with sampled gradients}
\label{ssec:ebm_mle}

The gradient of the negative log-likelihood under \cref{eq:ebm_definition} is
\begin{equation}
  -\nabla_\theta \log \ptheta(x)
  \;=\; \nabla_\theta \Etheta(x)
     \;-\; \EE_{x' \sim \ptheta}\!\left[\nabla_\theta \Etheta(x')\right].
  \label{eq:ebm_mle_grad}
\end{equation}
The data term $\nabla_\theta \Etheta(x)$ is direct.
The model expectation is the part that requires samples from $\ptheta$, typically via Markov chain Monte Carlo~\citep{hinton2002training,carreira2005contrastive}.
Each gradient step incurs the cost of running a sampler until it has approximately mixed.
For energy parameterizations as expressive as a transformer, sampling does not mix quickly~\citep{nijkamp2019learning,du2019implicit}, and exact MLE becomes operationally a budgeted approximation rather than a literal objective.

\subsection{Contrastive divergence and its persistent variant}
\label{ssec:cd}

\citet{hinton2002training} introduced contrastive divergence (CD-$k$), which replaces the long MCMC run in \cref{eq:ebm_mle_grad} with a short chain of $k$ steps starting from a data point.
The resulting estimator is biased but cheap.
Persistent contrastive divergence~\citep{tieleman2008training} maintains chains across parameter updates so that the negative samples track a slowly-changing $\ptheta$.
\citet{du2019implicit} and \citet{nijkamp2019learning} apply these to deep image energies with sampler buffers.
The empirical pattern is that CD trains and stabilizes only when the energy landscape is regular enough for short chains to be informative;
for very expressive energies, the bias of CD becomes the limiting factor and persistent variants are typically used.

\subsection{Noise-contrastive estimation}
\label{ssec:nce}

Noise-contrastive estimation~\citep{gutmann2010noise,gutmann2012noise} converts density estimation into binary classification: distinguish data samples from samples drawn from a known noise distribution $p_n$.
The unnormalized energy enters the classification logit, and $\log Z_\theta$ becomes a learnable scalar parameter trained jointly.
Asymptotically, the classifier learns $\log(\poftxt / p_n)$, and the recovered density is correct up to the noise distribution's coverage of the data~\citep{ma2018noise}.
The cost moves from sampling $\ptheta$ to choosing $p_n$.
A noise distribution that misses regions of the data renders the classifier blind to them~\citep{gutmann2014likelihood}.
Conditional NCE~\citep{mnih2012fast} adapts the method to language modeling by sampling negatives from a unigram base.

\subsection{Score matching}
\label{ssec:score_matching}

Score matching~\citep{hyvarinen2005score} avoids $\Ztheta$ entirely by matching $\nabla_x \log \ptheta(x)$ to $\nabla_x \log \poftxt(x)$ in expectation.
Because the gradient of $\log \Ztheta$ with respect to $x$ is zero, the score depends only on $\Etheta$:
\begin{equation}
  \nabla_x \log \ptheta(x) \;=\; -\nabla_x \Etheta(x).
\end{equation}
The original Fisher-divergence loss requires the trace of the Hessian of $\Etheta$, which is expensive for deep networks.
Two reformulations sidestep this.
\paragraph{Denoising score matching.}
\citet{vincent2011connection} prove that score matching against the smoothed density $\widetilde{p}(\tilde{x}) = \int q_\sigma(\tilde{x} \mid x) \poftxt(x) \mathrm{d}x$ for a known noise kernel $q_\sigma$ reduces to predicting the noise $\nabla_{\tilde{x}} \log q_\sigma(\tilde{x} \mid x)$.
This is the loss now standard in diffusion models~\citep{song2019generative,ho2020denoising}.
\paragraph{Sliced score matching.}
\citet{song2020sliced} replace the Hessian trace with random projections, giving an unbiased estimator with a single backward pass per sample.

Score matching has been adapted to text via discrete-variable extensions and to mixed sequence-embedding settings~\citep{meng2022concrete,lou2023discrete}.
The price of avoiding $\Ztheta$ at training is that the model is implicit in the score, and sampling requires running an iterative procedure (\cref{ssec:langevin}).

\subsection{Ratio matching and discrete energies}
\label{ssec:ratio_matching}

For discrete variables, the derivative-based score-matching identities do not apply.
\citet{hyvarinen2007some} introduce ratio matching: for each datum and each coordinate, the model matches the ratio of $\ptheta$ between flipping and not flipping the coordinate.
Generalized score matching for discrete data~\citep{lyu2009interpretation,grathwohl2021ono} extends this to wider notions of local perturbation.
In the language modeling context, \citet{deng2020residual} and \citet{lou2023discrete} use related ideas to define energies whose training loss is a coordinate-wise comparison rather than a full normalization.

\subsection{Joint energy / classifier hybrids}
\label{ssec:jem}

\citet{grathwohl2019your} observe that a logit-producing classifier defines an unconditional density via $\log p(x) = \lse(\mathrm{logits}(x))$ up to a constant.
Treating the classifier as both a discriminative model and an EBM gives the joint energy-based model (JEM).
Training combines the standard cross-entropy classification loss with a generative term computed via MCMC, sharing the partition function with the classification head.
The mechanism is consistent with the Hopfield reframe in \cref{ssec:hopfield}: a softmax-headed classifier is already an EBM.

\subsection{Stein and kernelized objectives}
\label{ssec:stein}

Kernelized Stein discrepancy~\citep{liu2016kernelized,chwialkowski2016kernel} provides an unnormalized goodness-of-fit objective that uses only $\nabla_x \log \ptheta(x)$.
Training an energy by minimizing the empirical Stein discrepancy avoids $\Ztheta$ at the cost of choosing a kernel and accepting bias proportional to the kernel's ability to detect mass discrepancies.

\subsection{What the training survey leaves}
\label{ssec:training_leaves}

Each method in \cref{tab:ebm_training} pays the cost of $\Ztheta$ somewhere.
Score-matching, denoising, and Stein methods push it into inference, where sampling becomes an iterative procedure (next section).
NCE pushes it into the choice of noise distribution and into the asymptotic bias if the noise misses modes.
MCMC-based estimators pay it directly in mixing time, which is the binding constraint on training large discrete energies.

The honest comparison with autoregressive language modeling is therefore not ``EBM versus AR.''
It is ``where does the cost of $\Ztheta$ go, and is the destination cheaper than $T$ local softmaxes plus the mode-covering objective?''
For domains with a hard verifier --- including formal-methods generation, code, structured outputs, and constraint satisfaction --- the cost of an inference-time iterative procedure can be amortized across a small number of verifier-bounded segments, which is precisely the setting in which the partition-function obstacle is most plausibly evaded in practice.
We return to this in \cref{sec:conclusion}.
